%!Mode:: "TeX:UTF-8"
%!TEX TS-program = xelatex
%arara: xelatex
\documentclass{ctexart}
\newif\ifpreface
\prefacetrue
\input{../../../global/all}
\begin{document}
\large
\setlength{\baselineskip}{1.2em}
\ifpreface
\input{../../../global/preface}
\newgeometry{left=2cm,right=2cm,top=2cm,bottom=2cm}
\else
\newgeometry{left=2cm,right=2cm,top=2cm,bottom=2cm}
\maketitle
\fi
%from_here_to_type
\begin{problem}\label{pro:1}
  假设 \(X = (X_t)_{t∈\mathbb{T}}\) 是 \((\Omega, \mathcal{F}, \mathbb{P})\) 上的随机过程, 状态空间为 \((E, \mathcal{E} )\). 证明下述结论:
  \begin{enumerate}
    \item \(X : (\Omega, \mathcal{F}) → (E^{\mathbb{T}}, \mathcal{E}^{\mathbb{T}})\), \(\omega \mapsto (X_t(\omega))_{ t \in \mathbb{T} }\) 是可测映射;
    \item \(X \)的有限维分布族唯一决定它的分布.
  \end{enumerate}
\end{problem}
\begin{solution}
  \begin{enumerate}
    \item \(\mathcal{E}^{\mathbb{T}} \) 可由所有有限维柱集生成，所以我们只需证明有限维柱集的原像是可测的。
      考查有限维柱集\(A:= \{f \in \mathcal{E}^{\mathbb{T}}: f(t_i)\in A_i,i=1,\cdots,n\} \)，其中\(t_{1},\cdots,t_n \in \mathbb{T},A_{1},\cdots,A_n \in \mathcal{E} \)。
      由定义\(X^{-1}(A) = \{\omega \in \Omega: X_{t_i}(\omega) \in A_i,i=1,\cdots,n\} = \bigcap_{i=1}^{n}X_{t_i}^{-1}(A_i) \)。
      由定义知\(X_{t_i}^{-1}(A_i) \)可测，故\(X^{-1}(A) \)可测。
    \item \(\mu_X \)是\(\mathcal{E}^{\mathbb{T}} \)上的测度，由于\(\mathcal{E}^{\mathbb{T}} \)可由有限维柱集生成，
      故由测度扩张的唯一性可知\(\mu_X \)由其在有限维柱集上的取值决定，从而\(X \)的有限维分友族唯一决定它的分布。
  \end{enumerate}
\end{solution}

\begin{problem}\label{pro:2}
  用下面的步骤完成讲义定理 11.1.1 的证明.
  其中状态空间为 \(E = \mathbb{R}\)，时间集为\( \mathbb{T} = [0, \infty)\)。
  \begin{enumerate}
    \item 假设 \(X = (X_t)_{t \geq 0} \)是一个中心化的 Gauss 过程, 即 \(X \)的任意有限维分布服从广义 Gauss 分布,
      且 \(\mathbb{E}X_t = 0, \forall t \geq 0\). 证明:\( X\) 的协方差函数\( K(s, t) := \mathbb{E}X_sX_t, s, t \geq 0 \)是一个对称，非负定函数。
      并且，\(X \)在有限时间集 \(I := \{t_1, t_2,\cdots, t_n\} \)上的分布\( \mu^X_I := \mathbb{P} \circ (X_{t_1}, \cdots , X_{t_n})^{-1}\)
      是 \(\mathbb{R}^I = \mathbb{R}^n\) 上的广义 Gauss 分布 \(N(0, A_I)\),
      其中 \(A_I = (K(t_i, t_j))_{1 \leq i,j \leq n}\)。
    \item
      假设 \(K \)是一个对称, 非负定函数. 对于任何 \(I := \{t_1,\cdots, t_n\}\),
      记 \(\mathbb{R}^n\) 上的广义 Gauss 分布 \(N(0, A_I)\) 为 \(\mu_I\), 其中 \(A_I\) 同上一问。
      再记 \(\mathcal{I} \)为 \(\mathbb{T} \)的有限子集全体. 证明:\( \{\mu_I : I \in \mathcal{I} \}\) 满足相容性 (consistent) 条件. 从而,
      利用 Kolmogorov 相容性定理得到, 存在 (分布意义下) 唯一的, 协方差函数为 \(K \)的中心化 Gauss 过程.
  \end{enumerate}
\end{problem}

\begin{solution}
  \begin{enumerate}
    \item 显然有
      \( K(s,t) = \mathbb{E}[X_sX_t] = \mathbb{E}[X_tX_s] = K(t,s), \)
      故 \(K\) 是对称函数。

      再任取 \(n \in \mathbb{N}\), \(t_1,\cdots,t_n \geq 0\) 及 \(c_1,\cdots,c_n \in \mathbb{R}\)，有
      \( \sum_{i,j=1}^n c_ic_jK(t_i,t_j) = \sum_{i,j=1}^n c_ic_j \mathbb{E}[X_{t_i}X_{t_j}] = \mathbb{E}\left[\left(\sum_{i=1}^n c_iX_{t_i}\right)^2\right] \geq 0. \)
      因此 \(K\) 是非负定函数。

      对于有限时间集
      \( I = \{t_1,\cdots,t_n\}, \)
      由于 \(X\) 是中心化 Gauss 过程，
      \( (X_{t_1},\cdots,X_{t_n}) \)
      服从某个广义 Gauss 分布。

      又因为
      \( \mathbb{E}[X_{t_i}] = 0, 1 \leq i \leq n, \)
      所以其均值向量为 \(0\)。

      此外，
      \( \mathrm{Cov}(X_{t_i},X_{t_j}) = \mathbb{E}[X_{t_i}X_{t_j}] = K(t_i,t_j), \)
      故其协方差矩阵为
      \( A_I = (K(t_i,t_j))_{1 \leq i,j \leq n}. \)

      从而
      \( \mu_I^X = \mathbb{P} \circ (X_{t_1}, \cdots , X_{t_n})^{-1} = N(0,A_I). \)

    \item 设 \(K\) 是一个对称非负定函数。对于任意有限集
      \( I = \{t_1,\cdots,t_n\}, \)
      定义
      \( A_I = (K(t_i,t_j))_{1 \leq i,j \leq n}. \)

      由于 \(K\) 非负定，对任意 \(c_1,\cdots,c_n \in \mathbb{R}\)，有
      \( \sum_{i,j=1}^n c_ic_jK(t_i,t_j) \geq 0, \)
      因此矩阵 \(A_I\) 半正定，从而广义 Gauss 分布
      \( \mu_I := N(0,A_I) \)
      是良定义的。

      下证 \(\{\mu_I : I \in \mathcal{I}\}\) 满足相容性条件。

      设
      \( I = \{t_1,\cdots,t_n\}, J = \{t_1,\cdots,t_n,t_{n+1},\cdots,t_m\}, \)
      且 \(I \subset J\)。

      记自然投影
      \( \pi_{JI} : \mathbb{R}^J \to \mathbb{R}^I \)
      为
      \( \pi_{JI}(x_1,\cdots,x_m) = (x_1,\cdots,x_n). \)

      若
      \( Y = (Y_{t_1},\cdots,Y_{t_m}) \sim \mu_J = N(0,A_J), \)
      则由于 Gauss 向量在线性映射下仍为 Gauss 向量，
      \( \pi_{JI}(Y) = (Y_{t_1},\cdots,Y_{t_n}) \)
      仍为中心化 Gauss 向量。

      其协方差矩阵为
      \( \mathrm{Cov}(Y_{t_i},Y_{t_j}) = K(t_i,t_j), 1 \leq i,j \leq n, \)
      即恰为 \(A_I\)。

      因此
      \( \mu_J \circ \pi_{JI}^{-1} = N(0,A_I) = \mu_I. \)
      故 \(\{\mu_I : I \in \mathcal{I}\}\) 满足相容性条件。

      由 Kolmogorov 相容性定理，存在随机过程
      \( X = (X_t)_{t \geq 0}, \)
      使得对任意有限集
      \( I = \{t_1,\cdots,t_n\}, \)
      都有
      \( (X_{t_1},\cdots,X_{t_n}) \sim \mu_I = N(0,A_I). \)

      因此 \(X\) 的任意有限维分布都是广义 Gauss 分布，且均值为 \(0\)，故 \(X\) 是中心化 Gauss 过程。

      又对任意 \(s,t \geq 0\)，有
      \( \mathbb{E}[X_sX_t] = K(s,t), \)
      故 \(X\) 的协方差函数为 \(K\)。

      最后，由于过程的有限维分布完全由 \(\{\mu_I\}\) 决定，因此该过程在分布意义下唯一。
  \end{enumerate}
\end{solution}
\begin{problem}
  考虑离散状态空间 \((E,\mathcal{E})\)，时间集 \(\mathbb{T} = \mathbb{N}\)。记
  \( \Omega := E^{\mathbb{N}}, \mathcal{F} := \mathcal{E}^{\mathbb{N}}, X_n(\omega) := \omega(n), \quad \forall \omega \in \Omega,\ n \geq 0. \)

  令
  \( p : E \times E \to [0,1] \)
  满足
  \( \sum_{y \in E} p(x,y) = 1, \forall x \in E. \)

  任给 \(E\) 上的分布
  \( \nu = (\nu_x)_{x \in E}, \)
  对于 \(n \geq 0\) 和时间集
  \( I_n = \{0,1,\cdots,n\}, \)
  定义 \(E^{I_n} = E^{n+1}\) 上的分布 \(\mu_{I_n}\) 为
  \( \mu_{I_n}(\{(x_0,\cdots,x_n)\}) = \nu_{x_0}p(x_0,x_1)\cdots p(x_{n-1},x_n), \forall x_i \in E,\ 0 \leq i \leq n. \)

  证明如下结论：
  \begin{enumerate}
    \item \(\{\mu_{I_n} : n \geq 0\}\) 可以扩张成一族满足相容性条件的有限维分布族。
    \item 用 Kolmogorov 相容性定理和典则实现证明：存在 \((\Omega,\mathcal{F})\) 上唯一的概率测度，记作 \(\mathbb{P}^{\nu}\)，使得 \((\Omega,\mathcal{F},\mathbb{P}^{\nu})\) 上的随机过程
      \( X = (X_n)_{n \geq 0} \)
      的有限维分布族与第一问中的有限维分布族一致。
    \item 对于 \(x \in E\)，当 \(\nu = \delta_x\) 时，令
      \( \mathbb{P}^x := \mathbb{P}^{\delta_x}. \)
      证明
      \( \mathbb{P}^{\nu} = \sum_{x \in E} \nu_x \mathbb{P}^x. \)
  \end{enumerate}
\end{problem}

\begin{solution}
  \begin{enumerate}
    \item 对于任意有限时间集
      \( I = \{t_0,t_1,\cdots,t_n\}, 0 \leq t_0 < t_1 < \cdots < t_n, \)
      定义 \(E^I\) 上的概率测度
      \( \mu_I(\{(x_{t_0},\cdots,x_{t_n})\}) := \nu_{x_{t_0}} p(x_{t_0},x_{t_1}) \cdots p(x_{t_{n-1}},x_{t_n}). \)

      当 \(I = I_n = \{0,1,\cdots,n\}\) 时，该定义与题目中的 \(\mu_{I_n}\) 一致，因此
      \( \{\mu_{I_n} : n \geq 0\} \)
      被扩张为一族有限维分布族。

      下证其满足相容性条件。

      设
      \( I = \{t_0,\cdots,t_n\} \subset J = \{s_0,\cdots,s_m\}, \)
      并记
      \( \pi_{JI} : E^J \to E^I \)
      为自然投影。

      任取
      \( (x_{t_0},\cdots,x_{t_n}) \in E^I, \)
      则
      \( \pi_{JI}^{-1}(\{(x_{t_0},\cdots,x_{t_n})\}) \)
      由所有在 \(I\) 上取值固定，而在 \(J \setminus I\) 上任意取值的点组成。

      因此
      \( \mu_J \circ \pi_{JI}^{-1}(\{(x_{t_0},\cdots,x_{t_n})\}) \)
      等于对所有补充坐标求和。利用
      \( \sum_{y \in E} p(x,y) = 1, \)
      可逐步消去所有额外变量，最终得到
      \( \mu_J \circ \pi_{JI}^{-1}(\{(x_{t_0},\cdots,x_{t_n})\}) = \nu_{x_{t_0}} p(x_{t_0},x_{t_1}) \cdots p(x_{t_{n-1}},x_{t_n}). \)

      即
      \( \mu_J \circ \pi_{JI}^{-1} = \mu_I. \)
      故该有限维分布族满足相容性条件。

    \item 由第一问知，
      \( \{\mu_I : I \subset \mathbb{N},\ I \text{ 有限}\} \)
      满足 Kolmogorov 相容性条件。

      因此由 Kolmogorov 相容性定理，存在概率空间
      \( (\Omega',\mathcal{F}',\mathbb{P}^{\nu}) \)
      及随机过程
      \( Y = (Y_n)_{n \geq 0}, \)
      使得对任意有限集
      \( I = \{t_0,\cdots,t_n\}, \)
      有
      \( (Y_{t_0},\cdots,Y_{t_n}) \sim \mu_I. \)

      现在取
      \( \Omega = E^{\mathbb{N}}, \mathcal{F} = \mathcal{E}^{\mathbb{N}}, \)
      并定义典则过程
      \( X_n(\omega) := \omega(n). \)

      Kolmogorov 定理进一步说明：存在唯一概率测度
      \( \mathbb{P}^{\nu} \)
      在 \((\Omega,\mathcal{F})\) 上，使得对任意有限集
      \( I = \{t_0,\cdots,t_n\}, \)
      及任意
      \( (x_{t_0},\cdots,x_{t_n}) \in E^I, \)
      有
      \( \mathbb{P}^{\nu} \bigl( X_{t_0}=x_{t_0}, \cdots, X_{t_n}=x_{t_n} \bigr) = \mu_I(\{(x_{t_0},\cdots,x_{t_n})\}). \)

      即典则过程
      \( X = (X_n)_{n \geq 0} \)
      的有限维分布族恰为第一问构造的有限维分布族。

    \item 令
      \( \mathbb{P}' := \sum_{x \in E} \nu_x \mathbb{P}^x. \)

      下证
      \( \mathbb{P}' = \mathbb{P}^{\nu}. \)

      只需证明 \(\mathbb{P}'\) 与 \(\mathbb{P}^{\nu}\) 有相同的有限维分布。

      任取
      \( I_n = \{0,1,\cdots,n\} \)
      及
      \( (x_0,\cdots,x_n) \in E^{n+1}, \)
      有
      \begin{align*}
        &\mathbb{P}' (X_0=x_0,\cdots,X_n=x_n) \\
        ={}&
        \sum_{x \in E}
        \nu_x
        \mathbb{P}^x
        (X_0=x_0,\cdots,X_n=x_n).
      \end{align*}

      由于在 \(\mathbb{P}^x\) 下，
      \( X_0 = x \)
      几乎处处成立，因此只有 \(x=x_0\) 时对应项非零，从而
      \begin{align*}
        &\mathbb{P}'
        (X_0=x_0,\cdots,X_n=x_n)
        \\
        ={}&
        \nu_{x_0}
        \mathbb{P}^{x_0}
        (X_1=x_1,\cdots,X_n=x_n).
      \end{align*}

      再由 \(\mathbb{P}^{x_0}\) 的有限维分布定义，
      \( \mathbb{P}^{x_0} (X_1=x_1,\cdots,X_n=x_n) = p(x_0,x_1)\cdots p(x_{n-1},x_n). \)

      因此
      \( \mathbb{P}' (X_0=x_0,\cdots,X_n=x_n) = \nu_{x_0} p(x_0,x_1)\cdots p(x_{n-1},x_n). \)

      这正是 \(\mu_{I_n}\) 的定义，即
      \( \mathbb{P}' (X_0=x_0,\cdots,X_n=x_n) = \mathbb{P}^{\nu} (X_0=x_0,\cdots,X_n=x_n). \)

      因而 \(\mathbb{P}'\) 与 \(\mathbb{P}^{\nu}\) 有相同的有限维分布。

      由第二问中的唯一性，
      \( \mathbb{P}^{\nu} = \mathbb{P}' = \sum_{x \in E} \nu_x \mathbb{P}^x. \)
  \end{enumerate}
\end{solution}
\begin{problem}
  仍然考虑离散状态空间 \((E,\mathcal{E})\)，时间集 \(\mathbb{T} = \mathbb{N}\)。令
  \( p : E \times E \to [0,1] \)
  满足
  \( \sum_{y \in E} p(x,y) = 1, \forall x \in E. \)

  假设
  \( X = (X_n)_{n \geq 0} \)
  是概率空间 \((\Omega,\mathcal{F},\mathbb{P})\) 上的随机过程，满足：对于任何 \(n \geq 0\),
  \( \mathbb{P}(X_{n+1}=x_{n+1}\mid X_n,\cdots,X_0) = p(X_n,x_{n+1}), \forall x_{n+1} \in E. \)

  记 \(X_0\) 的分布为
  \( \nu_x := \mathbb{P}(X_0=x), \forall x \in E. \)

  证明下述结论：
  \begin{enumerate}
    \item \(X\) 作为 \((\Omega,\mathcal{F},\mathbb{P})\) 上的随机过程服从 \((\nu,p)\)-分布。
    \item 任取并固定满足 \(\nu_x > 0\) 的 \(x \in E\)。定义
      \( \mathbb{Q}(A) := \mathbb{P}(A \mid X_0=x) = \frac{\mathbb{P}(A \cap \{X_0=x\})}{\mathbb{P}(X_0=x)}, \forall A \in \mathcal{F}. \)
      证明：\(\mathbb{Q}\) 是 \((\Omega,\mathcal{F})\) 上的概率测度。
    \item \(X=(X_n)_{n \geq 0}\) 是 \((\Omega,\mathcal{F},\mathbb{Q})\) 上的随机过程，它服从 \((\delta_x,p)\)-分布。
  \end{enumerate}
\end{problem}

\begin{solution}
  \begin{enumerate}
    \item 只需证明 \(X\) 的有限维分布与上一题构造的 \((\nu,p)\)-分布一致。

      任取
      \( n \geq 0, x_0,\cdots,x_n \in E. \)

      由条件概率公式，
      \begin{align*}
        &\mathbb{P}(X_0=x_0,\cdots,X_n=x_n)
        \\
        ={}&
        \mathbb{P}(X_n=x_n \mid X_0=x_0,\cdots,X_{n-1}=x_{n-1})
        \\
        & \times
        \mathbb{P}(X_0=x_0,\cdots,X_{n-1}=x_{n-1}).
      \end{align*}

      由题设中的条件，
      \( \mathbb{P}(X_n=x_n \mid X_0=x_0,\cdots,X_{n-1}=x_{n-1}) = p(x_{n-1},x_n). \)

      因此
      \begin{align*}
        &\mathbb{P}(X_0=x_0,\cdots,X_n=x_n)
        \\
        ={}&
        p(x_{n-1},x_n)
        \mathbb{P}(X_0=x_0,\cdots,X_{n-1}=x_{n-1}).
      \end{align*}

      反复迭代可得
      \begin{align*}
        &\mathbb{P}(X_0=x_0,\cdots,X_n=x_n)
        \\
        ={}&
        \mathbb{P}(X_0=x_0)
        p(x_0,x_1)\cdots p(x_{n-1},x_n)
        \\
        ={}&
        \nu_{x_0}
        p(x_0,x_1)\cdots p(x_{n-1},x_n).
      \end{align*}

      这正是上一题中 \((\nu,p)\)-分布对应的有限维分布，因此 \(X\) 服从 \((\nu,p)\)-分布。

    \item 由于 \(\nu_x > 0\)，故
      \( \mathbb{P}(X_0=x) > 0, \)
      从而 \(\mathbb{Q}\) 良定义。

      显然
      \( \mathbb{Q}(\Omega) = \frac{\mathbb{P}(\Omega \cap \{X_0=x\})}{\mathbb{P}(X_0=x)} = 1. \)

      又对任意 \(A \in \mathcal{F}\)，有
      \( \mathbb{Q}(A) = \frac{\mathbb{P}(A \cap \{X_0=x\})}{\mathbb{P}(X_0=x)} \geq 0. \)

      设 \((A_n)_{n \geq 1}\) 两两不交，则
      \begin{align*}
        \mathbb{Q}\left(\bigcup_{n=1}^{\infty} A_n\right)
        &=
        \frac{
          \mathbb{P}
          \left(
            \left(
              \bigcup_{n=1}^{\infty} A_n
            \right)
            \cap
            \{X_0=x\}
          \right)
        }{
          \mathbb{P}(X_0=x)
        }
        \\
        &=
        \frac{
          \mathbb{P}
          \left(
            \bigcup_{n=1}^{\infty}
            (A_n \cap \{X_0=x\})
          \right)
        }{
          \mathbb{P}(X_0=x)
        }
        \\
        &=
        \sum_{n=1}^{\infty}
        \frac{
          \mathbb{P}(A_n \cap \{X_0=x\})
        }{
          \mathbb{P}(X_0=x)
        }
        \\
        &=
        \sum_{n=1}^{\infty}
        \mathbb{Q}(A_n).
      \end{align*}

      因此 \(\mathbb{Q}\) 是 \((\Omega,\mathcal{F})\) 上的概率测度。

    \item 下证 \(X\) 在 \((\Omega,\mathcal{F},\mathbb{Q})\) 下服从 \((\delta_x,p)\)-分布。

      任取
      \( n \geq 0, x_0,\cdots,x_n \in E. \)

      由 \(\mathbb{Q}\) 的定义，
      \begin{align*}
        &\mathbb{Q}(X_0=x_0,\cdots,X_n=x_n)
        \\
        ={}&
        \frac{
          \mathbb{P}
          (
            X_0=x_0,\cdots,X_n=x_n,X_0=x
          )
        }{
          \mathbb{P}(X_0=x)
        }.
      \end{align*}

      当 \(x_0 \neq x\) 时，上式显然为 \(0\)。

      当 \(x_0=x\) 时，由第一问，
      \begin{align*}
        &\mathbb{Q}(X_0=x,\cdots,X_n=x_n)
        \\
        ={}&
        \frac{
          \mathbb{P}(X_0=x,\cdots,X_n=x_n)
        }{
          \mathbb{P}(X_0=x)
        }
        \\
        ={}&
        \frac{
          \nu_x
          p(x,x_1)\cdots p(x_{n-1},x_n)
        }{
          \nu_x
        }
        \\
        ={}&
        p(x,x_1)\cdots p(x_{n-1},x_n).
      \end{align*}

      综上，
      \begin{align*}
        &\mathbb{Q}(X_0=x_0,\cdots,X_n=x_n)
        \\
        ={}&
        \delta_x(x_0)
        p(x_0,x_1)\cdots p(x_{n-1},x_n).
      \end{align*}

      这正是 \((\delta_x,p)\)-分布对应的有限维分布。

      因此 \(X\) 在 \((\Omega,\mathcal{F},\mathbb{Q})\) 下服从 \((\delta_x,p)\)-分布。
  \end{enumerate}
\end{solution}
\end{document}
